% ===========================================================================
% Statistical Consistency — Per-Pixel Significance Testing
% ===========================================================================
% Requires: amsmath, amssymb, booktabs, siunitx, graphicx
% Figures : porous_512x512x128spp_difference_map.png
%           porous_512x512x128spp_z_diagnostic.png
% ===========================================================================

\subsection{Per-Pixel Statistical Consistency}
\label{sec:stat-consistency}

To rigorously confirm that the GPU wavefront solver produces
statistically indistinguishable results from the CPU depth-first
reference, we perform a per-pixel hypothesis testing analysis on a
shared radiative rendering scene.  The analysis goes beyond simple
visual comparison: it subjects every pixel to a formal significance
test, then applies multiple-comparison corrections to control the
false-discovery rate across the entire image.

%--- 5.X.1  Test Design ---------------------------------------------------
\subsubsection{Choice of indicators}

A Monte Carlo estimator at pixel $(x,y)$ yields an expected value
$\hat{E}$ and a standard error $\mathrm{SE} = \sqrt{V/N}$, where $V$ is
the sample variance and $N$ the number of independent realisations.
Because both the GPU and CPU solvers report $(\hat{E},\mathrm{SE})$
per pixel, we can construct a two-sample $Z$-statistic without
requiring access to the underlying path-level samples:
%
\begin{equation}\label{eq:z-pixel}
  Z_{x,y}
  = \frac{\bigl|\hat{E}_{\mathrm{gpu}}(x,y)
          - \hat{E}_{\mathrm{cpu}}(x,y)\bigr|}
         {\sqrt{\mathrm{SE}_{\mathrm{gpu}}^{2}(x,y)
              + \mathrm{SE}_{\mathrm{cpu}}^{2}(x,y)}}.
\end{equation}
%
Under the null hypothesis that both solvers sample the same underlying
distribution, $Z_{x,y}$ follows a half-normal distribution $|{\mathcal{N}(0,1)}|$
with mean $\sqrt{2/\pi}\approx 0.7979$ and standard deviation
$\sqrt{1-2/\pi}\approx 0.6028$.  Any systematic implementation
divergence would shift the $Z$~distribution toward higher values.

Three complementary indicators are reported:

\begin{enumerate}
\item \textbf{Uncorrected rejection rate.}
      The fraction of pixels satisfying $Z_{x,y}>1.96$ (two-sided
      $\alpha=0.05$).  Under the null this fraction converges to
      exactly~5\%; a value materially above~5\% signals the presence
      of systematic bias.

\item \textbf{Bonferroni correction.}
      The per-pixel significance level is set to
      $\alpha_{B}=\alpha/N_{\mathrm{valid}}$, yielding a
      family-wise error rate (FWER) $\le\alpha$.  This is the most
      conservative test: \emph{any} pixel surviving the Bonferroni
      threshold represents a near-certain localised discrepancy.

\item \textbf{Benjamini--Hochberg (BH) false-discovery-rate (FDR)
      control} \cite{benjamini1995}.
      The sorted $p$-values $p_{(1)}\le\cdots\le p_{(N)}$ are compared
      against the linear threshold $p_{(k)}\le (k/N)\,\alpha$.
      This controls the expected proportion of false discoveries at
      $\le\alpha$ while retaining substantially higher statistical
      power than Bonferroni, making it the primary acceptance criterion.
\end{enumerate}

The rationale for using all three is \emph{triangulation}: the raw rate
tests calibration of the $Z$~statistic itself; Bonferroni provides a
conservative upper bound on localised errors; and BH--FDR gives the
most powerful practical verdict.  When these three indicators are
mutually consistent, the conclusion is robust against the choice of
correction method.

A fourth diagnostic---the Kolmogorov--Smirnov (KS) test of $Z_{x,y}$
against the theoretical half-normal---assesses global distributional
agreement.  Because the KS test becomes extremely powerful at large~$N$
(rejecting even negligible deviations), we interpret the KS statistic
$D_n$ rather than its $p$-value: $D_n < 0.03$ indicates negligible
departure from the theoretical distribution.

%--- 5.X.2  Experimental setup --------------------------------------------
\subsubsection{Experimental setup}

The validation scene is a porous aluminium foam with a backing plate,
bounded by convective and fixed-temperature conditions
(Section~\ref{sec:scene}), rendered at $512\times512$ pixels with
$N_{\mathrm{spp}}=128$ samples per pixel on both solvers.
At this sample count the per-pixel standard error is
$\overline{\mathrm{SE}}\approx\SI{2.62}{\kelvin}$ (temperature range
\SIrange{680}{850}{\kelvin}), providing adequate statistical power to
detect biases of order $\SI{1}{\kelvin}$.

Pixels with $\mathrm{SE}_{\mathrm{combined}}<10^{-12}$ or zero
expected value on both sides (background sky rays that do not
intersect geometry) are excluded, leaving
$N_{\mathrm{valid}}=\num{153976}$ out of \num{262144} total pixels
(58.7\%).

%--- 5.X.3  Results --------------------------------------------------------
\subsubsection{Results}

\paragraph{Mean difference map.}
Figure~\ref{fig:diff-map} presents the temperature fields, the signed
pixel-wise difference $\Delta T = \hat{E}_{\mathrm{gpu}} -
\hat{E}_{\mathrm{cpu}}$, and the $Z$-statistic map.  The mean
difference across all valid pixels is
$\overline{\Delta T}=\SI{+0.006}{\kelvin}$, three orders of magnitude
below the per-pixel standard error, indicating \emph{no detectable
systematic bias}.  The standard deviation of $\Delta T$ is
$\SI{4.41}{\kelvin}$, consistent with the expected combined
noise $\sqrt{2}\,\overline{\mathrm{SE}}\approx\SI{3.70}{\kelvin}$
(the modest excess arises from the heavy-tailed SE distribution at
geometric edges).  The maximum absolute difference is
$\SI{34.2}{\kelvin}$ at a single edge pixel, which nevertheless yields
$Z=4.55$---below the Bonferroni threshold.

\paragraph{$Z$-distribution diagnostics.}
Table~\ref{tab:z-stats} and Figure~\ref{fig:z-diag} summarise the
$Z$-distribution.  The observed mean $\bar{Z}=0.806$ deviates from
the theoretical value $\sqrt{2/\pi}=0.798$ by only 1.0\%, and the
standard deviation is $0.606$ versus $0.603$ expected.
The Q--Q plot (Figure~\ref{fig:z-diag}, right panel) closely follows
the identity line up to $Z\approx3.5$, with a marginal heavy tail
beyond---attributable to the geometric edges where the CLT convergence
at $N=128$ is slower.  The KS statistic is $D_n=0.0226$, confirming
that the observed distribution is practically half-normal despite the
formal $p$-value being zero (an artefact of
$N_{\mathrm{valid}}=\num{153976}$).

\begin{table}[t]
  \centering
  \caption{Per-pixel $Z$-statistic summary and multiple-comparison
           results for the porous scene
           ($512\times512$, 128\,spp).}
  \label{tab:z-stats}
  \sisetup{round-mode=places}
  \begin{tabular}{@{}lS[round-precision=4]@{}}
    \toprule
    \textbf{$Z$-distribution} & \textbf{Value} \\
    \midrule
    Mean $Z$            & 0.8057 \\
    Median $Z$          & 0.6847 \\
    Std.\ dev.\ $Z$    & 0.6056 \\
    Max $Z$             & 4.5502 \\
    \midrule
    Theoretical mean $\sqrt{2/\pi}$ & 0.7979 \\
    Deviation from theory & {1.0\,\%} \\
    \bottomrule
  \end{tabular}
  \hspace{1em}
  \begin{tabular}{@{}lrr@{}}
    \toprule
    \textbf{Correction} & \textbf{Rejects} & \textbf{Rate} \\
    \midrule
    None ($Z>1.96$)          &  7\,858 & 5.10\,\% \\
    Bonferroni ($Z>5.11$)    &      0  & 0.000\,\% \\
    BH--FDR ($\alpha=0.05$)  &      0  & 0.00\,\% \\
    \bottomrule
  \end{tabular}
\end{table}

\paragraph{Multiple-comparison verdict.}
The uncorrected rejection rate of 5.10\,\% is statistically
indistinguishable from the nominal 5\,\% expected under pure Monte
Carlo noise, confirming that the $Z$-test is well-calibrated.
Under the Bonferroni correction ($\alpha_B =
\num{3.25e-7}$, threshold $Z>5.11$), \emph{zero} pixels are
rejected---meaning no single pixel exhibits a discrepancy that
could not be explained by chance alone in a family of
\num{153976}~simultaneous tests.  The BH--FDR procedure likewise
rejects \emph{zero} pixels, confirming the absence of any cluster of
localised biases that might survive a less conservative correction.

\begin{figure}[t]
  \centering
  \includegraphics[width=\linewidth]{porous_512x512x128spp_difference_map}
  \caption{Per-pixel comparison of the porous scene
           ($512\times512$, 128\,spp).
           \textit{Top row}: temperature maps from the CPU (left) and
           GPU (right) solvers.
           \textit{Bottom left}: signed difference
           $\Delta T = \hat{E}_{\mathrm{gpu}}-\hat{E}_{\mathrm{cpu}}$
           (colour centred at zero).
           \textit{Bottom right}: $Z$-statistic map.  No contours
           appear because both Bonferroni and FDR reject sets are
           empty.}
  \label{fig:diff-map}
\end{figure}

\begin{figure}[t]
  \centering
  \includegraphics[width=\linewidth]{porous_512x512x128spp_z_diagnostic}
  \caption{Diagnostic of the per-pixel $Z$-distribution.
           \textit{Left}: histogram of observed $Z$ (blue) versus the
           theoretical folded-normal density $2\varphi(z)$ (dashed).
           Vertical lines mark $z=1.96$ and the Bonferroni threshold
           $z=5.11$.
           \textit{Right}: Q--Q plot against the half-normal
           distribution.  Points follow the identity line closely,
           with a marginal heavy tail above $Z\approx3.5$.}
  \label{fig:z-diag}
\end{figure}

%--- 5.X.4  Interpretation ------------------------------------------------
\subsubsection{Interpretation}

Three independent lines of evidence converge to the same conclusion:

\begin{enumerate}
\item The \textbf{global bias} is negligible:
      $\overline{\Delta T}=\SI{+0.006}{\kelvin}$ on a
      \SIrange{680}{850}{\kelvin} temperature field, corresponding to
      a relative bias of $3.5\times10^{-5}$.

\item The \textbf{pixel-wise noise structure} is consistent with
      independent Monte Carlo sampling: the $Z$-distribution matches
      the half-normal to within 1\,\% in both mean and variance, and
      the uncorrected 5.10\,\% rejection rate agrees with the
      $\alpha=0.05$ type-I error rate.

\item \textbf{No localised discrepancy} survives multiple-comparison
      correction: both Bonferroni (FWER control) and BH--FDR yield
      zero rejections among \num{153976} simultaneous tests.
\end{enumerate}

These results demonstrate that the GPU wavefront implementation is a
statistically faithful reproduction of the CPU depth-first reference
at the per-pixel level.  All observed differences are fully accounted
for by the intrinsic variance of the Monte Carlo estimator, with no
evidence of algorithmic, numerical, or synchronisation-induced biases.
